% GAL1897 W03 -- Évariste Galois, Proof of a Theorem on Periodic
% Continued Fractions.
% Assigned source: physical PDF pages 030--037 / JP2 leaves L0029--L0036,
% printed pages 1--8. UTF-8 LaTeX fragment.
% Primary authority: bounded 600-ppi JP2 leaves, navigated against the
% canonical 96-page PDF. Printed errors are preserved by \GalPrintedError.
% See PAGE_DISPOSITIONS.csv, PRINTED_ERRORS.csv, and EVIDENCE_MANIFEST.csv.

\providecommand{\GalSourcePageStart}[4]{}%
\providecommand{\GalSourcePageBoundary}[4]{\space}%
\providecommand{\GalSourcePageBoundaryInWord}[4]{}%
\providecommand{\GalPrintedError}[2]{#1}%
\providecommand{\GalCopySpecificMarkImage}[4]{}%
\providecommand{\GalPrinterSignature}[1]{}%
\providecommand{\GalArticleEndRule}{\par\smallskip\begin{center}\rule{0.18\linewidth}{0.4pt}\end{center}}%
\providecommand{\GalHistoricalFootnote}[1]{\footnote{#1}}%

% Namespaced staircase-fraction machinery. It reproduces the left-to-right
% descending typography of the 1897 displays without rasterizing mathematics.
\makeatletter
\@ifundefined{GalStrut}{\newcommand{\GalStrut}[1][16pt]{\rule{0pt}{#1}}}{}%
\@ifundefined{GalCfrac}{\def\GalCfrac{\global\@tempdimb\z@\cfrac}}{}%
\@ifundefined{GalHlap}{\def\GalHlap#1{{\setbox\@tempboxa\hbox{#1}%
  \@tempdima=\wd\@tempboxa\global\advance\@tempdimb.45\@tempdima
  \hbox to.55\@tempdima{\unhbox\@tempboxa\hss}}}}{}%
\@ifundefined{GalCfracCorrect}{\def\GalCfracCorrect{\kern\@tempdimb\relax}}{}%
\makeatother

\GalSourcePageStart{030}{L0029}{1}{1}

% ENSEG:W03:PDF030:0001
\begin{center}
{\Large\bfseries MATHEMATICAL\par}
\vspace{0.8em}
{\small\bfseries WORKS\par}
\vspace{0.8em}
{\LARGE\bfseries OF ÉVARISTE GALOIS.\par}
\vspace{1.5em}
\rule{0.70\linewidth}{0.4pt}\par
\vspace{1.6em}
{\large\bfseries I. — ARTICLES PUBLISHED BY GALOIS.\par}
\vspace{1.6em}
{\bfseries PROOF\par}
\vspace{0.45em}
{\scriptsize\bfseries OF A\par}
\vspace{0.45em}
{\bfseries THEOREM ON PERIODIC CONTINUED FRACTIONS\textsuperscript{(1)}.\par}
\GalCopySpecificMarkImage{030}{L0029}{copy-specific non-letterpress mark beside article title}{figures/PDF030_L0029_TITLE_MARK_ENHANCED.png}%
\vspace{0.7em}
\rule{0.12\linewidth}{0.4pt}
\end{center}

% ENSEG:W03:PDF030:0002
\GalHistoricalFootnote{Volume XIX of M. Gergonne's \emph{Annales de Mathématiques},
page 294 (1828--1829). Galois was then a pupil at the Collège Louis-le-Grand.
\textsc{(J. Liouville.)}}
\GalPrinterSignature{E. G.}

% ENSEG:W03:PDF030:0003
It is known that if, by Lagrange's method, one expands one root of a quadratic
equation as a continued fraction, that continued fraction will be periodic.
The same is true of one root of an equation of arbitrary degree if that root
is a root of a rational quadratic factor of the left-hand side of the given
equation; in that case the equation will have at least one other root that is
also periodic. In either case the continued fraction may be immediately
periodic or not imme%
\GalSourcePageBoundaryInWord{031}{L0030}{2}{2}%
diately periodic; but in the latter case there will at least be one transformed
equation having a root that is immediately periodic.

% ENSEG:W03:PDF031:0004
Now, when an equation has two periodic roots belonging to the same rational
quadratic factor, and one of them is immediately periodic, there is a rather
singular relation between these two roots which appears not to have been
noticed before and which may be expressed by the following theorem.

\medskip
% ENSEG:W03:PDF031:0005
\noindent\textsc{Theorem.} — \emph{If one root of an equation of arbitrary
degree is an immediately periodic continued fraction, the equation must have
another root, likewise periodic, obtained by dividing negative unity by the
same periodic continued fraction written in reverse order.}
\medskip

% ENSEG:W03:PDF031:0006
\noindent\emph{Proof.} — To fix ideas, let us take periods of only four terms,
for the uniform course of the calculation shows that the same would hold if
we admitted a larger number. Let one root of an equation of arbitrary degree
be expressed as follows:
\[
x = a + \GalCfrac{1}{
      \,b + \GalHlap{$\cfrac{1}{
      \,c + \GalHlap{$\cfrac{1}{
      \,d + \GalHlap{$\cfrac{1}{
      \,a + \GalHlap{$\cfrac{1}{
      \,b + \GalHlap{$\cfrac{1}{
      \,c + \GalHlap{$\cfrac{1}{\GalStrut
      \,d + \raisebox{-2ex}{$\ddots$\rlap{\,;}} }$}}\;$}}\;$}}\;$}}\;$}}\;$}}\GalCfracCorrect
\]
the quadratic equation to which this root belongs, and which consequently
contains its conjugate, will be
\[
x = a + \GalCfrac{1}{
  \,b + \GalHlap{$\cfrac{1}{
  \,c + \GalHlap{$\cfrac{1}{
  \,d + \GalHlap{$\cfrac{1}{\quad x \quad } \;\rlap{;}$}}\;$}}\;$}}\GalCfracCorrect
\]
\GalSourcePageBoundary{032}{L0031}{3}{3}%
and from this we successively obtain
\begin{align*}
&a - x = - \GalCfrac{1}{
      \,b + \GalHlap{$\cfrac{1}{
      \,c + \GalHlap{$\cfrac{1}{
      \,d + \GalHlap{$\cfrac{1}{\quad x \quad } \;\rlap{,}$}}\;$}}\;$}}\GalCfracCorrect \qquad
&&\frac{1}{a-x} =
      - b + \GalCfrac{1}{
      \,c + \GalHlap{$\cfrac{1}{
      \,d + \GalHlap{$\cfrac{1}{\quad x \quad } \;\rlap{,}$}}\;$}}\GalCfracCorrect
\\[\baselineskip]
& b + \frac{1}{a-x} =
          - \GalCfrac{1}{
      \,c + \GalHlap{$\cfrac{1}{
      \,d + \GalHlap{$\cfrac{1}{\quad x \quad } \;\rlap{,}$}}\;$}}\GalCfracCorrect
&&\GalCfrac{1}{
      \,b + \GalHlap{$\cfrac{1}{
      \,a - x \, }\;$}}\GalCfracCorrect
=     - c + \GalCfrac{1}{
      \,d + \GalHlap{$\cfrac{1}{\quad x \quad } \;\rlap{,}$}}\GalCfracCorrect
\\[\baselineskip]
&c + \GalCfrac{1}{
      \,b + \GalHlap{$\cfrac{1}{
      \,a - x \, }\;$}}\GalCfracCorrect
=         - \GalCfrac{1}{
      \,d + \GalHlap{$\cfrac{1}{\quad x \quad } \;\rlap{,}$}}\GalCfracCorrect
&&\GalCfrac{1}{
      \,c + \GalHlap{$\cfrac{1}{
      \,b + \GalHlap{$\cfrac{1}{
      \,a - x \, }\;$}}\;$}}\GalCfracCorrect
=     -d + \frac{1}{x} \, ,
\\[\baselineskip]
&       d + \GalCfrac{1}{
      \,c + \GalHlap{$\cfrac{1}{
      \,b + \GalHlap{$\cfrac{1}{
      \,a - x \, }$}}\;$}}\GalCfracCorrect
=         - \frac{1}{x} \,;
&&\GalCfrac{1}{
      \,d + \GalHlap{$\cfrac{1}{
      \,c + \GalHlap{$\cfrac{1}{
      \,b + \GalHlap{$\cfrac{1}{
      \,a - x \, }$}}\;$}}\;$}}\GalCfracCorrect
=     -x\, ,
\end{align*}
that is,
\[
  x =     - \GalCfrac{1}{
      \,d + \GalHlap{$\cfrac{1}{
      \,c + \GalHlap{$\cfrac{1}{
      \,b + \GalHlap{$\cfrac{1}{
      \,a - x\, } \;\rlap{;}$}}\;$}}$}}\GalCfracCorrect
\]
this is therefore the quadratic equation that gives the two roots in question.
But by repeatedly substituting for $x$ on its right-hand side that same
right-hand side, which is indeed its value, it gives
\[
  x =     - \GalCfrac{1}{
      \,d + \GalHlap{$\cfrac{1}{
      \,c + \GalHlap{$\cfrac{1}{
      \,b + \GalHlap{$\cfrac{1}{
      \,a + \GalHlap{$\cfrac{1}{
      \,d + \GalHlap{$\cfrac{1}{
      \,c + \GalHlap{$\cfrac{1}{
      \,b + \GalHlap{$\cfrac{1}{
      \,a + \raisebox{-2ex}{$\ddots$\rlap{\,;}} }\;$}}\;$}}\;$}}\;$}}\;$}}\;$}}\;$}}
\GalCfracCorrect
\]
this is therefore the other value of $x$ given by the equation; as one sees,
it equals $-1$ divided by the first value written in reverse order.

\GalSourcePageBoundary{033}{L0032}{4}{4}%
% ENSEG:W03:PDF033:0007
In the preceding argument we supposed that the given root was greater than
unity; but if we had
\[
  x =       \GalCfrac{1}{
      \,a + \GalHlap{$\cfrac{1}{
      \,b + \GalHlap{$\cfrac{1}{
      \,c + \GalHlap{$\cfrac{1}{
      \,d + \GalHlap{$\cfrac{1}{
      \,a + \GalHlap{$\cfrac{1}{
      \,b + \GalHlap{$\cfrac{1}{
      \,c + \GalHlap{$\cfrac{1}{
      \,d + \raisebox{-2ex}{$\ddots$\rlap{\,;}} }$}}\;$}}\;$}}\;$}}\;$}}\;$}}\;$}}
\GalCfracCorrect
\]
we would conclude, for one value of $\dfrac{1}{x}$,
\[
  \frac{1}{x} =   a + \GalCfrac{1}{
      \,b + \GalHlap{$\cfrac{1}{
      \,c + \GalHlap{$\cfrac{1}{
      \,d + \GalHlap{$\cfrac{1}{
      \,a + \GalHlap{$\cfrac{1}{
      \,b + \GalHlap{$\cfrac{1}{
      \,c + \GalHlap{$\cfrac{1}{
      \,d + \raisebox{-2ex}{$\ddots$\rlap{\,;}} }$}}\;$}}\;$}}\;$}}\;$}}\;$}}
\GalCfracCorrect
\]
the other value of $\dfrac{1}{x}$ would therefore, by the preceding argument, be
\[
  \frac{1}{x} =   - \GalCfrac{1}{
      \,d + \GalHlap{$\cfrac{1}{
      \,c + \GalHlap{$\cfrac{1}{
      \,b + \GalHlap{$\cfrac{1}{
      \,a + \GalHlap{$\cfrac{1}{
      \,d + \GalHlap{$\cfrac{1}{
      \,c + \GalHlap{$\cfrac{1}{
      \,b + \GalHlap{$\cfrac{1}{
      \,a + \raisebox{-2ex}{$\ddots$\rlap{\,,}} }$}}\;$}}\;$}}\;$}}\;$}}\;$}}\;$}}
\GalCfracCorrect
\]
whence, for the other value of $x$, we would obtain
\[
  x = - d + \GalCfrac{1}{
      \,c + \GalHlap{$\cfrac{1}{
      \,b + \GalHlap{$\cfrac{1}{
      \,a + \GalHlap{$\cfrac{1}{
      \,d + \GalHlap{$\cfrac{1}{
      \,c + \GalHlap{$\cfrac{1}{
      \,b + \GalHlap{$\cfrac{1}{
      \,a + \raisebox{-2ex}{$\ddots$} }$}}\;$}}\;$}}\;$}}\;$}}\;$}}
\GalCfracCorrect
\]
\GalSourcePageBoundary{034}{L0033}{5}{5}%
or
\[
  x = - 1 : \GalCfrac{1}{
      \,d + \GalHlap{$\cfrac{1}{
      \,c + \GalHlap{$\cfrac{1}{
      \,b + \GalHlap{$\cfrac{1}{
      \,a + \GalHlap{$\cfrac{1}{
      \,d + \GalHlap{$\cfrac{1}{
      \,c + \GalHlap{$\cfrac{1}{
      \,b + \GalHlap{$\cfrac{1}{
      \,a + \raisebox{-2ex}{$\ddots$\rlap{\,;}}}$}}\;$}}\;$}}\;$}}\;$}}\;$}}\;$}}
\GalCfracCorrect
\]
which falls exactly under our theorem.

% ENSEG:W03:PDF034:0008
Let $A$ be any immediately periodic continued fraction, and let $B$ be the
continued fraction obtained by reversing its period. If one root of an equation
is $x=A$, it must have another root $x=-\dfrac{1}{B}$. Now, if $A$ is positive
and greater than unity, $-\dfrac{1}{B}$ will be negative and lie between $0$
and $-1$; conversely, if $A$ is negative and lies between $0$ and $-1$,
$-\dfrac{1}{B}$ will be positive and greater than unity. Thus, when one root
of a quadratic equation is an immediately periodic continued fraction greater
than unity, the other must lie between $0$ and $-1$; conversely, if one lies
between $0$ and $-1$, the other must be positive and greater than unity.

% ENSEG:W03:PDF034:0009
Conversely, one can prove that if one of the two roots of a quadratic equation
is positive and greater than unity while the other lies between $0$ and $-1$,
then those roots can be expressed as immediately periodic continued fractions.
Indeed, continue to let $A$ be any positive immediately periodic continued
fraction greater than unity, and let $B$ be the immediately periodic continued
fraction obtained by reversing its period; like $A$, it too will be positive
and greater than unity. The first root of the given equation cannot have the form
\[
  x = p + \frac{1}{A},
\]
\GalSourcePageBoundary{035}{L0034}{6}{6}%
for then, by our theorem, the second would have to be
\[
  x = \GalPrintedError{a}{W03-PE001} + \frac{1}{-\dfrac{1}{B}}
    = \GalPrintedError{a}{W03-PE001} - B ;
\]
now $\GalPrintedError{a}{W03-PE001}-B$ could not lie between $0$ and
$\GalPrintedError{1}{W03-PE002}$ unless the integer part of $B$ were equal to
$p$, in which case the first value would be immediately periodic. Nor could
the first value of $x$ be $x=p+\dfrac{1}{q+\dfrac{1}{A}}$, for then the other
would be
\[
  x = p + \frac{1}{q - B} \qquad \text{or} \qquad
  x = p - \frac{1}{B - q};
\]
now, for this value to lie between $0$ and $-1$, it would first be necessary
that $\dfrac{1}{B-q}$ equal $p$ plus a fraction. Hence $B-q$ would have to be
less than unity, requiring $B$ to equal $q$ plus a fraction. Thus $q$ and $p$
would have to equal, respectively, the first two terms of the period belonging
to $B$, or the last two terms of the period belonging to $A$. Contrary to the
hypothesis, the value $x=p+\dfrac{1}{q+\dfrac{1}{A}}$ would therefore be
immediately periodic. An analogous argument would prove that the periods
cannot be preceded by a larger number of terms not belonging to them.

\GalCriticalNote{W03-PE001}{The conjugate expression must use $p$, not $a$}{Conjugating
$x=p+1/A$ gives $p+1/(-1/B)=p-B$. The two printed occurrences of $a$ are
inconsistent with the variable already established, so both the conjugation
and the interval argument must use $p-B$. No earlier explicit notice of this
letter error was found in the sources examined.}
\GalCriticalNote{W03-PE002}{The interval lies below zero}{If $B=p+r$ with
$0<r<1$, then $p-B=-r$, so $-1<p-B<0$. The printed interval ``between $0$
and $1$'' has the wrong sign; it should be between $0$ and $-1$. The converse
argument that follows already requires this corrected interval.}

% ENSEG:W03:PDF035:0010
Thus, when treating a numerical equation by Lagrange's method, we may be sure
that no periodic roots are to be expected until we encounter a transformed
equation having at least one positive root greater than unity and another
lying between $0$ and $-1$. If the root being pursued is indeed periodic,
its period will begin no later than at this transformed equation.

% ENSEG:W03:PDF035:0011
If one root of a quadratic equation is not only immediately periodic but also
symmetric---that is, if terms equally distant from the ends of the period are
equal---then $B=A$; consequently the two roots
\GalSourcePageBoundary{036}{L0035}{7}{7}%
will be $A$ and $-\dfrac{1}{A}$, and the equation will therefore be
\[
Ax^{2} - (A^{2}-1)x - A = 0.
\]
Conversely, every quadratic equation of the form
\[
ax^{2} - bx - a = 0
\]
has roots that are both immediately periodic and symmetric. Indeed, substituting
in turn infinity and $-1$ for $x$ gives positive results, whereas substituting
$x=1$ and $x=0$ gives negative results. We first see that the equation has one
positive root greater than unity and one negative root between $0$ and $-1$,
and hence that these roots are immediately periodic. Moreover, the equation
does not change when $x$ is replaced by $-\dfrac{1}{x}$. Thus, if $A$ is one
of its roots, the other is $-\dfrac{1}{A}$, so in this case $B=A$.

% ENSEG:W03:PDF036:0012
Let us apply these general results to the quadratic equation
\[
3x^{2} - 16x + 18 = 0;
\]
we first find that it has a positive root between $3$ and $4$. Setting
\[
x = 3 + \frac{1}{y},
\]
we obtain the transformed equation
\[
3y^{2} - 2y - 3 = 0,
\]
whose form tells us that the values of $y$ are both immediately periodic and
symmetric. Indeed, setting successively
\[
y = 1 + \frac{1}{z}, \qquad
z = 2 + \frac{1}{t}, \qquad
t = 1 + \frac{1}{u},
\]
we obtain the transformed equations
\[
2z^{2} - 4z - 3 = 0, \qquad
3t^{2} - 4t - 2 = 0, \qquad
3u^{2} - 2u - 3 = 0.
\]
The identity of the equations in $u$ and $y$ proves that the positive value
\GalSourcePageBoundary{037}{L0036}{8}{8}%
of $y$ is
\[
  y =   1 + \GalCfrac{1}{
      \,2 + \GalHlap{$\cfrac{1}{
      \,1 + \GalHlap{$\cfrac{1}{
      \,1 + \GalHlap{$\cfrac{1}{
      \,2 + \GalHlap{$\cfrac{1}{
      \,1 + \raisebox{-2ex}{$\ddots$\rlap{\,;}} }$}}\;$}}\;$}}\;$}}
\GalCfracCorrect
\]
its negative value will therefore be
\[
  y =     - \GalCfrac{1}{
      \,1 + \GalHlap{$\cfrac{1}{
      \,2 + \GalHlap{$\cfrac{1}{
      \,1 + \GalHlap{$\cfrac{1}{
      \,1 + \GalHlap{$\cfrac{1}{
      \,2 + \GalHlap{$\cfrac{1}{
      \,1 + \raisebox{-2ex}{$\ddots$\rlap{\,;}} }$}}\;$}}\;$}}\;$}}\;$}}
\GalCfracCorrect
\]
the two values of $x$ will therefore be
\[
\begin{aligned}
  x &= 3 + \GalCfrac{1}{
      \,1 + \GalHlap{$\cfrac{1}{
      \,2 + \GalHlap{$\cfrac{1}{
      \,1 + \GalHlap{$\cfrac{1}{
      \,1 + \GalHlap{$\cfrac{1}{
      \,2 + \GalHlap{$\cfrac{1}{
      \,1 + \raisebox{-2ex}{$\ddots$\rlap{\,,}} }$}}\;$}}\;$}}\;$}}\;$}}
\GalCfracCorrect
\qquad
  x &= \GalPrintedError{3 - \GalCfrac{1}{
      \,1 + \GalHlap{$\cfrac{1}{
      \,2 + \GalHlap{$\cfrac{1}{
      \,1 + \GalHlap{$\cfrac{1}{
      \,1 + \GalHlap{$\cfrac{1}{
      \,2 + \GalHlap{$\cfrac{1}{
      \,1 + \raisebox{-2ex}{$\ddots$\rlap{\,,}} }$}}\;$}}\;$}}\;$}}\;$}}
\GalCfracCorrect}{W03-PE003}
\end{aligned}
\]
the latter of which, by the familiar formula
\[
  p - \frac{1}{q}
= p - 1 + \GalCfrac{1}{
      \,1 + \GalHlap{$\cfrac{1}{q-1}\, ,$}}
\GalCfracCorrect
\]
becomes
\[
  x =   1 + \GalCfrac{1}{
      \,1 + \GalHlap{$\cfrac{1}{
      \,1 + \GalHlap{$\cfrac{1}{
      \,1 + \GalHlap{$\cfrac{1}{
      \,1 + \GalHlap{$\cfrac{1}{
      \,2 + \GalHlap{$\cfrac{1}{
      \,1 + \GalHlap{$\cfrac{1}{
      \,1 + \GalHlap{$\cfrac{1}{
      \,2 + \GalHlap{$\cfrac{1}{
      \,1 + \raisebox{-2ex}{$\ddots$\rlap{\ .}} }$}}\;$}}\;$}}\;$}}\;$}}\;$}}\;$}}\;$}}
\GalCfracCorrect
\]
\GalCriticalNote{W03-PE003}{The second root is $3-A$}{The printed staircase
expression $3-1/A$ is not a root of $3x^{2}-16x+18=0$. Direct substitution
shows that the correct second root is $3-A$, which is also the value consistent
with the terminal continued fraction. No earlier explicit notice of this
formula error was found in the sources examined.}
\GalArticleEndRule
